\documentclass[11pt,a4paper]{article}

\usepackage[spanish,es-noquoting]{babel}
\usepackage{fontspec}
\usepackage[margin=2.5cm]{geometry}
\usepackage{booktabs}
\usepackage{tabularx}
\usepackage{enumitem}
\usepackage{xcolor}
\usepackage[hidelinks]{hyperref}

\setlist{nosep, leftmargin=*}
\setlength{\parskip}{0.5em}
\setlength{\parindent}{0pt}

\newcommand{\decision}[1]{\textcolor{teal}{\textbf{Decisión:}} #1}
\newcommand{\abierta}[1]{\textcolor{orange!80!black}{\textbf{Abierta:}} #1}

\title{Propuesta de frontend\\[0.3em]
  \large Calendario armonizado de focos de calor MODIS--VIIRS}
\author{Equipo Hots-Spots --- NASA Space Apps Challenge 2026}
\date{14 de noviembre de 2026 \quad·\quad Borrador 0.1}

\begin{document}
\maketitle

\section{Propósito de este documento}

Plantear \emph{qué} va a hacer la aplicación web, \emph{para quién} y \emph{cómo} la
vamos a construir, antes de escribir código. La armonización de los datos la lleva
otra parte del equipo; este plan fija qué necesita el frontend de esos datos (sección
\ref{sec:contrato}) para que los dos trabajos avancen en paralelo.

\section{Objetivo}

Que una persona no experta, en menos de un minuto, pueda responder:
\begin{itemize}
  \item ¿Este mes es peor de lo normal en mi zona?
  \item ¿Cuál fue el peor año, y en qué meses se concentra la actividad?
  \item ¿Cuántos focos hubo en esta área y periodo? (y descargarlos)
\end{itemize}
con una serie 2000--2025 en la que un foco significa lo mismo sin importar el satélite.

\subsection*{Usuarios}

\begin{tabularx}{\linewidth}{@{}lX@{}}
\toprule
Usuario & Qué necesita \\
\midrule
Gestores de emergencias y bomberos & Respuesta rápida: ¿está siendo un mes anómalo? \\
Gestores del territorio & Comparar zonas y épocas; planificar temporadas \\
Científicos & Series consistentes, incertidumbre explícita, descarga en CSV \\
\bottomrule
\end{tabularx}

\section{Principios de diseño}

\begin{enumerate}
  \item \textbf{Una pregunta, una pantalla.} Una sola vista principal; nada de menús
        profundos ni pestañas anidadas.
  \item \textbf{El dato es el protagonista.} Fondo neutro, tipografía sobria, un único
        color de acento. El color se reserva para codificar cantidad de focos.
  \item \textbf{Rampa de color accesible.} Secuencial, perceptualmente uniforme y apta
        para daltonismo (familia \emph{viridis}/\emph{inferno} o ColorBrewer \emph{YlOrRd}).
        Para anomalías, una rampa divergente con el centro en ``normal''.
  \item \textbf{Honestidad con la incertidumbre.} Los años 2000--2002 (solo Terra) y
        los meses con pocos focos se marcan como menos precisos, sin esconderlos.
  \item \textbf{Instantáneo.} Datos precalculados; nada de consultas a FIRMS en vivo.
  \item \textbf{Responsive y legible en móvil}, porque parte del público está en campo.
\end{enumerate}

\section{Funcionalidades del MVP}

\begin{tabularx}{\linewidth}{@{}clX@{}}
\toprule
\# & Funcionalidad & Detalle \\
\midrule
F1 & Selección de área & Dibujar un rectángulo/polígono en el mapa, o introducir
     latitud/longitud y radio. Opción de ``región completa''. \\
F2 & Selección de periodo & Rango de años y/o meses; atajo para un día concreto. \\
F3 & Calendario años $\times$ meses & Matriz de 26 filas (2000--2025) por 12 columnas,
     cada celda coloreada por focos armonizados. Es la vista central. \\
F4 & Comparación con lo normal & Para el periodo elegido: mediana histórica, anomalía
     (\% sobre la mediana) y posición en el ranking. Alternar calendario en
     \emph{valores absolutos} o \emph{anomalía}. \\
F5 & Rankings & Los 5 meses y los 5 años con más actividad del área. \\
F6 & Detalle al pasar/tocar & Celda $\to$ focos armonizados, focos brutos, sensor de
     origen, banda de incertidumbre. \\
F7 & Descarga & CSV con los focos de la búsqueda (con su peso) y CSV/PNG del calendario. \\
\bottomrule
\end{tabularx}

\subsection*{Fuera del MVP}
Cobertura global, datos NRT en tiempo real, predicción, cuentas de usuario.

\section{Bocetos de la interfaz}

Se harán en una fase posterior (sección \ref{sec:fases}). Estructura prevista de
la pantalla única, de arriba abajo:

\begin{enumerate}
  \item \textbf{Barra superior mínima:} nombre, selector de región, botón de ayuda.
  \item \textbf{Panel de consulta:} mapa pequeño con la herramienta de dibujo y el
        selector de periodo. En escritorio, a la izquierda; en móvil, plegable arriba.
  \item \textbf{Respuesta en una frase:} p.\,ej. ``Agosto de 2025: 38\,\% por encima
        de lo normal --- el 3.\textsuperscript{er} agosto más activo desde 2000''.
  \item \textbf{Calendario} años $\times$ meses con su leyenda.
  \item \textbf{Rankings y descarga}, discretos, al final.
\end{enumerate}

\section{Contrato de datos}\label{sec:contrato}

Lo que el frontend espera de la armonización. Es la pieza que hay que acordar primero
con el equipo de datos.

\begin{tabularx}{\linewidth}{@{}lX@{}}
\toprule
Producto & Contenido \\
\midrule
\texttt{celdas\_mes.parquet} & Agregado por celda de rejilla (p.\,ej. $0{,}1^\circ$) y mes:
  \texttt{lat}, \texttt{lon}, \texttt{mes}, \texttt{focos\_armonizados}
  (=\,$\sum$\texttt{peso}), \texttt{focos\_brutos}, \texttt{fuente}
  (\texttt{TERRA}/\texttt{MODIS}/\texttt{VIIRS}). Alimenta el calendario de cualquier área. \\
\texttt{k\_mes.json} & Constantes por mes del año y su rango intercuartil, para la banda
  de incertidumbre. \\
\texttt{focos/*.parquet} & Focos individuales con su \texttt{peso}, particionados por año,
  solo para la descarga en CSV. \\
\texttt{metadatos.json} & Región, periodo cubierto, filtros aplicados, versión del método. \\
\bottomrule
\end{tabularx}

\abierta{tamaño de celda de la rejilla ($0{,}1^\circ$ frente a $0{,}25^\circ$): equilibrio
entre precisión del área dibujada y peso de los archivos.}

\section{Arquitectura técnica}

\begin{tabularx}{\linewidth}{@{}lX@{}}
\toprule
Pieza & Propuesta \\
\midrule
Base & React 19 + TypeScript + Vite + Tailwind 4 (ya en el repo). \\
Mapa & MapLibre GL JS con un mapa base claro y sobrio, más una herramienta de dibujo. \\
Consultas & DuckDB-WASM leyendo los Parquet en el navegador: filtra por área y periodo
  en milisegundos, sin servidor de base de datos. \\
Gráficos & SVG propio para el calendario (es una cuadrícula simple); evaluar
  Observable Plot si hace falta más. \\
Despliegue & Sitio estático (GitHub Pages o similar) con los Parquet como archivos. \\
\bottomrule
\end{tabularx}

\decision{sin backend en el MVP. Todo es estático, lo que cumple ``respuesta instantánea''
y simplifica el despliegue.}

\abierta{si el volumen de focos individuales es demasiado para el navegador, la descarga
de CSV pasaría a una pequeña función en servidor.}

\section{Fases}\label{sec:fases}

\begin{tabularx}{\linewidth}{@{}clX@{}}
\toprule
Fase & Nombre & Entregable en \texttt{Documentacion/frontend/} \\
\midrule
0 & Plan & Este documento. \\
1 & Contrato de datos & Esquema acordado con el equipo de datos y un conjunto de
  datos de ejemplo (puede ser sintético) para no bloquearnos. \\
2 & Bocetos & Wireframes de escritorio y móvil, y la guía visual: paleta, tipografía,
  rampa de color y espaciado. \\
3 & Prototipo & Calendario y selector de periodo funcionando con los datos de ejemplo. \\
4 & Mapa y área & Dibujo de áreas y consulta con DuckDB-WASM. \\
5 & Comparación y descarga & Anomalías, rankings, CSV y PNG. \\
6 & Pulido & Accesibilidad (contraste, teclado, lectores de pantalla), rendimiento,
  textos de ayuda y prueba con usuarios no expertos (la prueba del minuto). \\
\bottomrule
\end{tabularx}

\subsection*{Prototipo del mapa (adelantado)}

Se adelantó parte de la fase 4 para tener algo visible pronto: un mapa (MapLibre GL,
mapa base \emph{Positron} de OpenFreeMap, sin clave) que pinta los focos armonizados
de India en una rejilla de $0{,}1^\circ$, con selector de año (2000--2025) y de mes.
Al pasar el ratón sobre una celda se ven los focos armonizados y las detecciones brutas;
un recuadro da el total del periodo y avisa cuando la fuente es solo Terra.

\begin{itemize}
  \item Datos: \texttt{analisis/05\_exportar\_mapa.sql} (\texttt{mise run datos:mapa})
        escribe un JSON por año en \texttt{frontend/public/datos/} (16\,MB en total,
        $\sim$0,6\,MB por año). Es un adelanto del contrato \texttt{celdas\_mes}, en JSON
        en vez de Parquet para no depender aún de DuckDB-WASM.
  \item Rampa YlOrRd en 6 escalones casi logarítmicos (\texttt{<3}, 3, 10, 30, 100, 300+
        focos por celda), la misma para mes y año completo.
\end{itemize}

\abierta{¿el mapa pasa a ser la vista principal, con el calendario como panel que se
actualiza según el área elegida en el mapa?}

\section{Criterios de aceptación}

\begin{itemize}
  \item Un usuario nuevo responde ``¿este mes es peor de lo normal?'' en menos de un minuto.
  \item Cualquier consulta de área y periodo responde en menos de 1\,s.
  \item El calendario no muestra salto en 2012 y marca 2000--2002 como menos precisos.
  \item La rampa de color se distingue con los tres tipos comunes de daltonismo.
  \item El CSV descargado coincide con lo que muestra el calendario.
\end{itemize}

\section{Preguntas abiertas}

\begin{enumerate}
  \item ¿``Normal'' es la mediana o el promedio? ¿De todo 2000--2025 o de una ventana
        móvil, por ejemplo los 10 años anteriores?
  \item ¿Qué región muestra la demo final: India (los datos que ya tenemos) o la región
        de la CONAE (Argentina)?
  \item ¿Idioma de la interfaz: solo español, o español e inglés para el jurado?
  \item ¿Mostrar el número absoluto de focos armonizados, o un índice relativo para no
        dar una falsa sensación de precisión?
\end{enumerate}

\end{document}
